The literature reviewed in Chapter \ref{ch:literature_review} collectively demonstrates that compact radar hardware is now sufficiently mature for UAV DAA applications. SWaP-compliant systems have been built, detection at DAA-relevant ranges has been achieved, and processing feasibility on constrained embedded platforms has been established. However, three important gaps are identified.

First, existing hardware demonstrations validate detection performance against 
large, high radar cross-section targets such as buildings and fixed-wing aircraft. 
No work in the UAV DAA literature characterises how compact radar systems perform 
against small moving airborne targets under realistic encounter conditions, leaving 
the practical sensitivity floor of these systems unknown.

Second, while GPU acceleration has been shown to offer substantial speedups for 
radar signal processing in non-UAV contexts, no work has benchmarked a 
GPU-based radar pipeline in a UAV DAA setting. Existing UAV DAA processing 
studies are limited to FPGA and ARM CPU implementations, evaluate single 
algorithm sequences on single hardware configurations, and do not profile 
architectural bottlenecks such as host--device data transfer and thread 
synchronisation overhead. As a result, it remains unclear which algorithm 
combinations and implementation strategies are most effective when jointly 
optimising for sensitivity, latency, and throughput on modern embedded GPU 
platforms.

Third, sensitivity algorithms for low-SNR monopulse radar --- including methods 
for resolving unresolved targets within a single angular cell and for jointly 
treating detection and angle estimation across channels --- have only been 
validated under idealised white Gaussian noise conditions with single simulated 
targets. None of these methods have been integrated into a complete end-to-end 
range--Doppler tracking pipeline, and their behaviour under realistic 
multi-target, non-Gaussian noise conditions is unknown.

This thesis addresses all three gaps. Small moving airborne targets are 
included explicitly in both simulation and real data collection, and performance 
is measured not only by detection range but by false alarm behaviour, track 
continuity, and sensitivity in low-SNR regimes. A comprehensive selection of 
radar processing and tracking algorithms is benchmarked on embedded GPU 
hardware, with explicit profiling of host--device transfer costs, memory layout 
decisions, and synchronisation overhead across multiple hardware configurations. 
Finally, monopulse sensitivity algorithms are evaluated within a complete 
multi-target tracking pipeline under realistic noise conditions, with track 
quality measured alongside compute cost. Together, these experiments provide 
the first consolidated, systems-level evaluation of UAV DAA radar processing 
that jointly addresses sensitivity, latency, and throughput on modern embedded 
platforms.
